\hypertarget{forward}{}
\hypertarget{strictly-forward-validation}{%
\chapter{Strictly forward
validation}\label{strictly-forward-validation}}

This chapter designs an evaluation in which every model choice and estimate
used to form the ranking is fixed \emph{before} any future observation is
touched --- the discipline that separates description from prediction
(\S\ref{rq-estimand}).

\hypertarget{timeline}{%
\section{Timeline}\label{timeline}}

At each cutoff $c$ the procedure is:
\begin{enumerate}
\tightlist
\item
  take the 120-month training window $\{c-119,\dots,c\}$;
\item
  fit the betas, the joint HAC covariance $V$, and the empirical-Bayes
  hyperparameters $\mu,\tau$ on that window;
\item
  freeze the posterior scores and the ranking;
\item
  observe the next 12 months $\{c+1,\dots,c+12\}$; and
\item
  evaluate each portfolio's realised future alpha using the \emph{frozen}
  training betas,
\end{enumerate}
\begin{equation}\label{eq:ch11-fwdalpha}
A_i^{+} = \frac{1}{12}\sum_{t=c+1}^{c+12}\bigl(y_{it}-\hat\beta_{i,\text{train}}^{\top} f_t\bigr).
\end{equation}
The beta is frozen deliberately: re-estimating it on the evaluation year
would let future data redefine the very target being evaluated.

\textbf{The formal meaning of ``look-ahead free.''} Let $\mathcal{F}_c$ be
the information available at the cutoff. The ranking $R_c$ must be
$\mathcal{F}_c$-measurable, while the outcome is realised strictly later:
\begin{equation}\label{eq:ch11-measurable}
R_c \in \mathcal{F}_c,
\qquad
A_c^{+} \in \mathcal{F}_{c+12}\setminus\mathcal{F}_c .
\end{equation}
Every estimated beta, covariance, hyperparameter, scaling rule,
tie-breaking rule and portfolio assignment must be determined by
$\mathcal{F}_c$; the leakage assertions below enforce exactly this in code.

\hypertarget{evaluation-statistics}{%
\section{Evaluation statistics}\label{evaluation-statistics}}

Two statistics summarise each evaluation window $w$. The first is the
Spearman rank correlation (\S\ref{prelim-rank}) between the 25 training
scores and the 25 future alphas $A_i^{+}$. The second is the average future
alpha of the training top quintile minus that of the training bottom
quintile,
\begin{equation}\label{eq:ch11-spread}
D_w = 12\times10{,}000\times\Bigl(\tfrac15\!\sum_{i\in\text{top}_w} A_i^{+}
 - \tfrac15\!\sum_{i\in\text{bottom}_w} A_i^{+}\Bigr)\quad\text{bps/year},
\end{equation}
where $\text{top}_w$ and $\text{bottom}_w$ are the five highest- and
lowest-ranked portfolios at cutoff $w$; averaging within quintiles reduces
dependence on any single extreme.

Correlations are bounded, so their sampling distribution is skewed near
$\pm1$ and they must not be averaged on the raw scale. Following
\S\ref{prelim-rank}, each annual correlation is put on the
variance-stabilising \emph{Fisher} scale, averaged there with HAC
inference, and mapped back:
\begin{equation}\label{eq:ch11-fisher}
z_w = \operatorname{atanh}(\rho_w) = \tfrac12\log\frac{1+\rho_w}{1-\rho_w},
\qquad \bar\rho = \tanh(\bar z).
\end{equation}
Because the $z_w$ are serially dependent across (annual, non-overlapping)
windows, the long-run variance of their mean is estimated with a four-lag
HAC (\S\ref{prelim-hac}); the 95\% interval on the $z$ scale is
$\bar z \pm 1.96\,\mathrm{SE}_{\mathrm{HAC}}(\bar z)$, with both endpoints
mapped back through $\tanh$. The spread $D_w$ is aggregated the same way,
without transformation. This annual mean-with-HAC construction is exactly
the equal-predictive-accuracy machinery formalised in Chapter~12.

\hypertarget{leakage-assertions}{%
\section{Leakage assertions}\label{leakage-assertions}}

\begin{verbatim}
assert train_dates.max() < eval_dates.min()
assert len(set(train_dates) & set(eval_dates)) == 0
assert beta_timestamp < eval_dates.min()
assert rank_timestamp < eval_dates.min()
assert hyperparameter_timestamp < eval_dates.min()
assert preprocessing_fit_timestamp < eval_dates.min()
\end{verbatim}

\textbf{Paper result.} Across 89 non-overlapping evaluation years, mean
rank correlation is 0.102 with HAC 95\% interval {[}0.055, 0.149{]}. The
top-minus-bottom quintile spread averages 243 bps/year, but its IQR of
498 bps exceeds its mean and negative years remain frequent.

\textbf{A positive average is not a timing rule.} Forward validation
shows association using past information. It does not yet prove that
annual re-estimation adds value over a simpler stable ranking.

